\documentclass[10pt]{article}
\providecommand{\FIGDIR}{fig}
\newcommand{\DIGESTDATE}{24 August 2026}
% ===================== PM Research Watch - shared preamble =====================
\usepackage[T1]{fontenc}
\usepackage[utf8]{inputenc}
\usepackage{amsmath,amssymb}
\usepackage{tgpagella}
\usepackage{tgheros}
\renewcommand{\sfdefault}{qhv}
\renewcommand{\ttdefault}{lmtt}
\usepackage{microtype}
\usepackage[english]{babel}

\usepackage[a4paper,top=17mm,bottom=17mm,left=17mm,right=17mm,
            headsep=6mm,footskip=10mm]{geometry}
\usepackage{graphicx}
\usepackage{xcolor}
\usepackage{colortbl}
\usepackage{tikz}
\usetikzlibrary{calc}
\usepackage[most]{tcolorbox}
\usepackage{titlesec}
\usepackage{fancyhdr}
\usepackage{booktabs}
\usepackage{array}
\usepackage{longtable}
\usepackage{enumitem}
\usepackage{multicol}
\usepackage{ragged2e}
\usepackage[hidelinks]{hyperref}

% ------------------------------------------------------------------ palette
\definecolor{Ink}    {HTML}{14262E}
\definecolor{Deep}   {HTML}{0F4C5C}
\definecolor{Teal}   {HTML}{1F7A8C}
\definecolor{Sky}    {HTML}{4FA3C4}
\definecolor{Sage}   {HTML}{6FA287}
\definecolor{Amber}  {HTML}{E3A008}
\definecolor{Clay}   {HTML}{B5643C}
\definecolor{Coral}  {HTML}{D1495B}
\definecolor{Violet} {HTML}{6A5B9E}
\definecolor{Slate}  {HTML}{4A6572}
\definecolor{Paper}  {HTML}{FBF9F5}
\definecolor{Mist}   {HTML}{EEEAE2}
\definecolor{Rule}   {HTML}{D9D3C9}

\pagecolor{Paper}
\color{Ink}

\hypersetup{colorlinks=true, linkcolor=Deep, urlcolor=Teal, citecolor=Teal,
            pdftitle={PM Research Watch}, pdfauthor={Automated daily digest}}

% ------------------------------------------------------------------ headings
\titleformat{\section}
  {\sffamily\bfseries\large\color{Deep}}{}{0pt}
  {}[{\vspace{1.2mm}\color{Amber}\titlerule[1.1pt]}]
\titlespacing*{\section}{0pt}{5.5mm}{2.8mm}

\titleformat{\subsection}
  {\sffamily\bfseries\normalsize\color{Teal}}{}{0pt}{}
\titlespacing*{\subsection}{0pt}{3.6mm}{1.4mm}

\setlist[itemize]{leftmargin=4.2mm, itemsep=1.1mm, topsep=1.3mm, parsep=0pt,
                  label=\textcolor{Teal}{\footnotesize$\blacktriangleright$}}

% ------------------------------------------------------------------ header/footer
\pagestyle{fancy}
\fancyhf{}
\renewcommand{\headrulewidth}{0pt}
\renewcommand{\footrulewidth}{0pt}
\fancyhead[L]{\sffamily\scriptsize\color{Slate}PM RESEARCH WATCH\;\textcolor{Rule}{\textbar}\;\DIGESTDATE}
\fancyhead[R]{\sffamily\scriptsize\color{Slate}Particulate matter monitoring \& health effects}
\fancyfoot[C]{\sffamily\scriptsize\color{Slate}\thepage}

% ------------------------------------------------------------------ components
% metric tile
\newcommand{\metric}[3]{%
  \begin{tcolorbox}[enhanced, colback=white, colframe=Rule, boxrule=0.5pt,
    arc=1.6mm, left=2.2mm, right=2.2mm, top=1.8mm, bottom=1.8mm,
    borderline west={1.7mm}{0pt}{#3}, nobeforeafter, valign=center,
    height=17mm, width=\linewidth]
    \RaggedRight
    {\sffamily\bfseries\fontsize{17}{17}\selectfont\color{#3}#1}\\[0.5mm]
    {\sffamily\fontsize{6.5}{7.4}\selectfont\color{Slate}#2}
  \end{tcolorbox}}

% pull-quote / key insight box
\newtcolorbox{keybox}[1][]{enhanced, colback=Mist, colframe=Mist,
  arc=1.6mm, left=3.2mm, right=3.2mm, top=2.4mm, bottom=2.4mm,
  borderline west={1.7mm}{0pt}{Amber}, #1}

% subtopic band
\newcommand{\band}[2]{%
  \vspace{3.2mm}\noindent
  \begin{tcolorbox}[enhanced, colback=#2, colframe=#2, arc=1.2mm,
    left=2.6mm, right=2.6mm, top=1.1mm, bottom=1.1mm, boxrule=0pt,
    width=\linewidth, nobeforeafter]
    {\sffamily\bfseries\small\color{white}#1}
  \end{tcolorbox}\vspace{1.6mm}}

% one paper entry: {short cite}{title}{journal}{doi}{body}
\newcommand{\paperentry}[5]{%
  \par\vspace{1.9mm}\noindent
  {\sffamily\bfseries\footnotesize\color{Ink}#1}\;%
  {\sffamily\fontsize{7}{8}\selectfont\color{Slate}#3}\\[0.4mm]
  {\footnotesize\itshape\color{Deep}#2}\\[0.6mm]
  {\footnotesize #4}\\[0.3mm]
  {\fontsize{6.6}{7.6}\selectfont\color{Slate}\texttt{doi:}\,\href{https://doi.org/#5}{\texttt{\detokenize{#5}}}}
  \par}

% caption style
\newcommand{\figcap}[1]{\par\vspace{1.0mm}
  {\fontsize{7.2}{8.6}\selectfont\color{Slate}\textsf{\textbf{Fig.}}\; #1\par}\vspace{1.5mm}}

\providecommand{\permil}{\textperthousand}

\setlength{\parindent}{0pt}
\setlength{\parskip}{0pt}
\renewcommand{\arraystretch}{1.18}


\begin{document}

% ==================================================================== MASTHEAD
\thispagestyle{empty}
\begin{tikzpicture}[remember picture, overlay]
  \fill[Deep] (current page.north west) rectangle
      ($(current page.north east)+(0,-67mm)$);
  \fill[Amber] ($(current page.north west)+(0,-67mm)$) rectangle
      ($(current page.north east)+(0,-68.4mm)$);
  \foreach \x/\y/\r in {22/12/0.9, 41/26/1.5, 63/9/0.7, 88/31/1.1, 112/17/1.9,
                        138/29/0.8, 159/13/1.3, 178/54/1.0, 196/21/1.6,
                        30/44/0.6, 74/57/1.2, 125/48/0.9, 168/6/0.8, 191/59/0.7,
                        52/18/1.0, 100/60/0.6, 148/20/1.1, 66/49/1.4, 205/45/0.8}
    \fill[Sky, opacity=0.30] ($(current page.north west)+(\x mm,-\y mm)$) circle (\r mm);
\end{tikzpicture}

\vspace*{4mm}
{\sffamily\fontsize{7.6}{9}\selectfont\color{Sky}%
 AUTOMATED DAILY LITERATURE DIGEST\;\textbullet\;PhD RESEARCH WATCH\par}
\vspace{2.2mm}
{\sffamily\bfseries\fontsize{27}{29}\selectfont\color{white}%
 Particulate Matter\par}
{\sffamily\fontsize{27}{29}\selectfont\color{Sky}%
 Monitoring \& Health Effects\par}
\vspace{5.0mm}
{\sffamily\fontsize{8.4}{10}\selectfont\color{white}%
 \mbox{Issue 2026-08-24}\;\;\textcolor{Amber}{\textbar}\;\;\mbox{Entry window 24 Aug 2026}\;\;
 \textcolor{Amber}{\textbar}\;\;\mbox{9 in scope, 21 screened out}\par}
\vspace{18mm}

% -------------------------------------------------------------- metric strip
\noindent
\begin{minipage}[t]{0.190\linewidth}\metric{9}{RECORDS IN SCOPE\\(21 SCREENED OUT)}{Deep}\end{minipage}\hfill
\begin{minipage}[t]{0.190\linewidth}\metric{4}{OF 10 CLUSTERS\\OCCUPIED}{Teal}\end{minipage}\hfill
\begin{minipage}[t]{0.190\linewidth}\metric{2}{ABSTRACT-FREE\\DEPOSITS (22\%)}{Sky}\end{minipage}\hfill
\begin{minipage}[t]{0.190\linewidth}\metric{3}{ATMOS. CHEM. PHYS.\\RECORDS (A FIRST)}{Amber}\end{minipage}\hfill
\begin{minipage}[t]{0.190\linewidth}\metric{2}{INTERVAL-BEARING\\EXPOSURE ESTIMATES}{Clay}\end{minipage}

\vspace{6mm}

% -------------------------------------------------------------- signal of day
\section*{\sffamily\bfseries\color{Deep}Signal of the day}
\vspace{-1.5mm}{\color{Amber}\rule{\linewidth}{1.1pt}}\vspace{2.5mm}

\begin{keybox}
{\sffamily\bfseries\small\color{Deep}The Crossref by-ISSN leg finally caught
\textit{Atmospheric Chemistry and Physics}, and the first four records it delivered
include the cleanest quantification this watch has seen of the error a light-scattering
sensor makes when it cannot tell a swollen aerosol particle from an activated droplet.}\\[1.4mm]
\footnotesize
Neuberger et al.\ (\textit{Atmos Chem Phys}) analysed the FAIRARI 2021/2022 fog campaign
at San Pietro Capofiume in the Po Valley with $\kappa$--K\"ohler theory and the MIMICA
large-eddy simulation. \textbf{Failing to account for hydrated, non-activated particles
inflated the inferred effective diameter by 81\,\% (11.6 to 21.0\,$\mu$m) and cut the
inferred cloud droplet number concentration by 87\,\% (97.4 to 12.4\,cm$^{-3}$)}, while
hydrated particles supplied \textbf{21\,\% of liquid water content} and accounted for
\textbf{36\,\% of sub-kilometre visibility events with no droplet activation at all}.
Median $\kappa$ separated cleanly by population --- fog residuals 0.45, interstitial
0.40, out-of-fog aerosol 0.34. This is a fog-microphysics paper, but the mechanism is
exactly the one that governs low-cost optical PM sensors at high relative humidity: the
instrument sizes a water-coated particle and reports the mass of a dry one. The
literature's usual answer is an empirical RH correction fitted to a co-location period,
which implicitly assumes a fixed $\kappa$; Neuberger et al.\ measure $\kappa$ varying by
a third \emph{within a single site and campaign} depending on which population the
particle belongs to. A humidity correction estimated during one aerosol regime should
therefore not be expected to transfer to another, and the residual after correction is
not noise --- it is composition.
\end{keybox}

\vspace{3mm}

% -------------------------------------------------------------- exec brief
\section*{\sffamily\bfseries\color{Deep}What changed today}
\vspace{-1.5mm}{\color{Amber}\rule{\linewidth}{1.1pt}}\vspace{2.5mm}

\footnotesize
\begin{itemize}
\item \textbf{A structural coverage gap closed, at least for one day.} The sweep caught
\textbf{four} \textit{Atmospheric Chemistry and Physics} deposits --- the first ACP
records it has ever returned --- of which \textbf{three shipped}; the fourth, a
Jungfraujoch radiocarbon CH$_4$/CO$_2$ study, is gas-phase only and was screened out.
ACP is one of the two journals the source audit named as
structurally invisible to PubMed, so until now the watch has been blind to it except
through Consensus, which is relevance-ranked rather than date-windowed and therefore
resurfaces old work. \textbf{This is a harvest-architecture result, not a literature
result}, and it should be read as such: the by-ISSN Crossref leg is now demonstrably
delivering the instrumentation-side coverage it was added for. Whether it does so
reliably needs several more days before it can be claimed.

\item \textbf{Consensus returned six on-topic sensor and exposure-modelling papers and
every one of them was a duplicate.} Low-cost sensor calibration, recalibration cadence,
calibration-strategy review, satellite--surface fusion in North-East India, an
urban--rural balanced estimation framework for China, and a WRF-Chem/satellite/ground
fusion scheme --- all six were already in \texttt{seen.json}, shipped between 26 July and
5 August. The dedupe key caught all six before a single word was written about them. This
is the expected behaviour of a relevance-ranked source against a date-windowed pipeline,
and it is worth stating plainly that on a thin day Consensus adds nothing to recall; its
value is in the topics PubMed cannot index, not in the recency of what it returns.

\item \textbf{Two mechanistic lung-cancer records from the same Chinese journal issue,
pointing at the same exposure from opposite ends.} Li \& Li review GPER as the receptor
through which PM$_{2.5}$'s estrogen-like organics and metal ions drive MAPK/ERK and
PI3K/AKT signalling; the Wumeng Mountain synthesis assembles the exposure--component--
effect--molecular chain for household smoky-coal lung cancer in the Yunnan--Guizhou
Plateau, where the disease concentrates in \textbf{never-smoking women} with a
\textbf{distinctive driver-gene spectrum} (low classical EGFR 19del/L858R, elevated
G719X, S768I, KRAS G12C, ALK and ROS1 fusions). Both are narrative reviews with no
pooled estimate and no formal search protocol reported in the deposited abstract, so
neither carries evidential weight on its own. What they carry jointly is a testable
claim: that a specific source-resolved PM composition maps to a specific somatic
mutation spectrum. That is the kind of claim source-apportionment-capable distributed
sensing could eventually address and mass-only networks cannot.

\item \textbf{The only intervals in the window are on self-reported exposure
categories.} Aulia Anwar et al.'s North Sumatra case-control (240/240) reports
\textbf{OR 2.53 (95\,\% CI 1.45--4.40)} for an unhealthy indoor environment and
\textbf{OR 1.91 (1.13--3.25)} for occupational exposure. Neither is a measured
particulate concentration --- both are questionnaire categories --- so the forest plot
names the exposure metric on its own axis rather than implying a PM contrast. The
same model's age ORs are demographic covariates and are excluded from the figure for
the reason recorded on 20 and 23 August.

\item \textbf{Two abstract-free Elsevier deposits again, both from \textit{Atmospheric
Environment}.} A machine-learning wind-correction comparison for coupled WRF--CMAQ
(dynamic assimilation vs static replacement) and a particulate PAH dry-deposition study
in the Bohai Sea Strait. Both ship as \textit{Metadata only}. Scholar Gateway was queried
for the wind-correction paper and cannot help --- its corpus ends May 2026 --- and
Consensus has not indexed either. The Elsevier abstract backlog first noted on 8 August
is therefore still open and now spans three weeks.
\end{itemize}

\vspace{3mm}

% ==================================================================== FIGURES
\section*{\sffamily\bfseries\color{Deep}Corpus at a glance}
\vspace{-1.5mm}{\color{Amber}\rule{\linewidth}{1.1pt}}\vspace{2.5mm}

\noindent\includegraphics[width=\linewidth]{\FIGDIR/f1_subtopics.png}
\figcap{Records per subtopic cluster. \textbf{Four of ten clusters are occupied} ---
exposure assessment \& modelling 3, sensing/instrumentation 2, mechanistic toxicology 2,
respiratory \& allergic 2. The six empty clusters are cardiovascular \& metabolic,
neuro / mental health, reproductive \& developmental, occupational \& indoor, burden /
policy \& mitigation, and other clinical endpoints. Yesterday's neuro pair and burden
record have no counterpart today; the balance has shifted from clinical epidemiology
toward atmospheric process and instrumentation, which is what the ACP arrivals do to a
nine-record day.}

\vspace{2mm}
\noindent\includegraphics[width=\linewidth]{\FIGDIR/f2_design_endpoint.png}
\figcap{Left --- study architecture, grouped. \textit{Modelling / inventory} 2,
\textit{Metadata only} 2, \textit{Review / synthesis} 2, and one each of measurement
campaign, experimental / toxicology and observational cross-sectional --- six
architectures over nine records, so the panel is close to flat by construction. Right ---
health endpoint. \textbf{Cancer 3} (both lung-cancer reviews plus the North Sumatra
case-control), respiratory 1, and \textbf{five of nine records carry no health endpoint
at all}: two instrument/atmospheric-process papers, one exposure-modelling paper, and the
two metadata-only deposits whose endpoint the publisher did not deposit. Those five are
split across two bars --- \textit{No health endpoint} 3 and \textit{Not extracted} 2 ---
rather than pooled, so the panel understates the no-endpoint share on a first reading.
Separately, the triage tiers are one A (Neuberger et al.), four B
and four C; the C block is the two abstract-free deposits plus the two records whose
exposure is a questionnaire category or an undeposited aerosol dose.}

\vspace{2mm}
\noindent\includegraphics[width=\linewidth]{\FIGDIR/f3_metric_geography.png}
\figcap{Left --- particle metric, as binned by the coarse grouper: \textit{PM
(unspeciated) / emissions} 6, \textit{PM$_{2.5}$ only} 2, \textit{size-resolved / number
distribution} 1. \textbf{The two records in the PM$_{2.5}$ bar do not measure PM$_{2.5}$
mass} --- they are binned there because their metric strings contain the token: Li \& Li
treat PM$_{2.5}$-borne estrogenic organics and metal ions as receptor ligands, and the
Wumeng synthesis names smoky-coal PM$_{2.5}$/black carbon inside a PAH-dominated mixture.
Read literally, \textbf{no record this window uses gravimetric PM$_{2.5}$ as its exposure
variable}; the metrics actually in play are hygroscopicity parameter and hydrated
diameter, nanoparticle growth rate, isoprene-SOA mass and subspecies composition,
particle-bound PAH, source-resolved coal-smoke composition, an inhaled AlCl$_3$ aerosol
and a self-reported household-environment category. Right --- geographic provenance.
\textbf{China 5}, Europe 1 (Po Valley), Southeast Asia 1 (North Sumatra), multi-region 2
--- and the multi-region bar is one genuinely pooled multi-continent growth-rate dataset
plus one metadata-only record whose study region the publisher did not deposit, not two
global studies.}

\vspace{2mm}
\noindent\includegraphics[width=\linewidth]{\FIGDIR/f5_forest.png}
\figcap{The window's only interval-bearing exposure estimates, both from the same
multivariate model in Aulia Anwar et al.\ (240 cases / 240 controls, North Sumatra).
\textbf{Both exposures are self-reported questionnaire categories, not measured
concentrations}, and the axis is labelled accordingly. Two points from one study is not
a synthesis and no pooled estimate is implied or drawn.}

\vspace{2mm}
\noindent\includegraphics[width=\linewidth]{\FIGDIR/f6_lifecourse.png}
\figcap{Life-course windows addressed, counted only where a record's study population is
explicitly anchored at that stage. \textbf{One of five windows occupied}, and it is
working-age adulthood: the North Sumatra case-control and the aluminium-ILD supportive
cohort. The perinatal, paediatric and adolescent windows are empty for the second
consecutive issue, and unlike yesterday the 65+ window is empty too --- no record this
window studies an older population.}

\vspace{2mm}
\noindent\includegraphics[width=\linewidth]{\FIGDIR/f4_heatmap.png}
\figcap{Subtopic $\times$ study-architecture cross-tabulation. \textbf{The only cell
holding more than one record is exposure assessment $\times$ \textit{Metadata only}} ---
i.e.\ the densest cell in the matrix is an artefact of two publisher deposits arriving
without abstracts, not a methodological concentration. Every other occupied cell holds
exactly one record, and the two narrative reviews sit in different rows (mechanistic
toxicology and respiratory) despite arriving from the same journal issue. With nine
records over six architectures the matrix is close to a permutation and reading a pattern
into it would be over-interpretation.}

\vspace{3mm}
% ==================================================================== DIGEST
\section*{\sffamily\bfseries\color{Deep}Paper-level digest}
\vspace{-1.5mm}{\color{Amber}\rule{\linewidth}{1.1pt}}\vspace{1mm}

{\fontsize{7.4}{9}\selectfont\color{Slate}Ordered by cluster. Each entry gives the
methodological core and the finding that matters, not an abstract paraphrase.\par}

\band{Sensing, forecasting \& instrumentation\hfill 2 records}{Amber}

\paperentry{Neuberger et al.\ 2026}
{Importance of hydrated aerosol particles for aerosol--fog relationships in the Italian Po Valley}
{Atmos Chem Phys}
{FAIRARI 2021/2022 aerosol--fog observations at San Pietro Capofiume, analysed with
$\kappa$--K\"ohler theory and the MIMICA large-eddy simulation. Median hygroscopicity
separated by population: \textbf{fog residuals $\kappa = 0.45$, interstitial particles
0.40, out-of-fog aerosol 0.34}, the gradient reflecting inorganic enrichment in activated
droplets. Excluding hydrated but non-activated particles from the analysis
\textbf{inflated effective diameter by 81\,\% (11.6 $\rightarrow$ 21.0\,$\mu$m) and
depressed droplet number concentration by 87\,\% (97.4 $\rightarrow$ 12.4\,cm$^{-3}$)};
hydrated particles supplied \textbf{21\,\% of liquid water content} and drove
\textbf{36\,\% of sub-kilometre visibility events without any activation}. The LES leg
shows fog prediction hinging on the largest dry particles. The limitation that binds is
site specificity --- one heavily polluted valley over two winters, with a $\kappa$
distribution shaped by Po Valley nitrate that will not transfer to a marine or biomass-
burning regime. For distributed low-cost sensing the transferable point is that the
hydrated/activated distinction is a \emph{composition} problem, so an RH correction
fitted at one site during one season inherits that site's $\kappa$ and no more.}
{10.5194/acp-26-11991-2026}

\paperentry{Zhang et al.\ 2026c}
{Buffering of atmospheric nanoparticle growth by temperature-dependent shifts in molecular composition, volatility and diffusivity}
{Atmos Chem Phys}
{A kinetic multiphase treatment applied to previously unreconciled field and chamber
nanoparticle growth-rate data. The standing puzzle is that observed growth rates are
smaller than theory predicts and \emph{less} sensitive to condensable vapour
concentration than theory predicts. The resolution offered is two antagonistic
temperature dependencies: \textbf{slow surface-to-bulk transport limits vapour uptake by
semi-solid, low-diffusivity particles, while Clausius--Clapeyron shifts in the volatility
distribution enhance growth at low temperature and low concentration}. Netted out, they
buffer the vapour-concentration dependence, which is why measured growth rates converge
on a few nm/h across widely varying conditions. No new measurement is reported --- this
is a reinterpretation of existing datasets, and its persuasiveness rests on the
volatility-distribution parameterisation rather than on independent validation. The
relevance here is upstream of sensing: the 50--100\,nm photochemical mode that Zhong
et al.\ flagged yesterday as invisible to mass-based assessment is produced by exactly
this growth process, and a mechanism that makes its rate insensitive to precursor
loading also makes it insensitive to the emission controls that reduce PM$_{2.5}$ mass.}
{10.5194/acp-26-12037-2026}

\band{Exposure assessment \& modelling\hfill 3 records}{Deep}

\paperentry{Zhang et al.\ 2026d}
{Improved representation of isoprene-derived secondary organic aerosol in CAM6-Chem reveals regional contrasts in its long-term changes over China}
{Atmos Chem Phys}
{An explicit isoprene-SOA formation scheme implemented in CAM6-Chem with high-NO$_x$
gas-phase epoxide precursors and branch-resolved heterogeneous uptake from both low- and
high-NO$_x$ pathways. Against ground observations the update \textbf{improves the SOA
normalised mean bias over China from $-76.7\,\%$ to $-51.6\,\%$} and reproduces the
compositional structure of four explicit ISOA subspecies. Over 2000--2019 the national
mean ISOA trend is weak because regional changes cancel: \textbf{the sharpest increase is
in Southwest China, where rising biogenic isoprene emission is the leading statistical
predictor at a 58.3\,\% regression-based contribution}, with the strongest decrease in
Shaanxi--Gansu--Ningxia. The honest reading of the headline number is that a $-51.6\,\%$
bias is still a factor-of-two underestimate --- the scheme narrows the gap without
closing it, and a regression-based attribution is not a causal decomposition. For
exposure work the consequence is directional: any health study using a CTM-derived
organic-aerosol surface over China is working with a field that is known-low by roughly
half, and non-uniformly so by region.}
{10.5194/acp-26-11967-2026}

\paperentry{Kim et al.\ 2026}
{Application of machine learning-based wind corrections to coupled WRF--CMAQ simulations: a comparison of dynamic assimilation and static replacement}
{Atmos Environ}
{\textbf{Metadata only.} Crossref by-ISSN sweep; title, authorship and DOI verified, no
abstract deposited. The comparison it names --- injecting an ML-corrected wind field into
CMAQ dynamically during the run versus statically replacing the WRF field beforehand ---
is a live methodological question for anyone producing exposure surfaces, because wind
error is the dominant transport error in urban CTM output and the two strategies have
very different mass-conservation properties. Nothing is asserted about which won. Scholar
Gateway cannot reach it (corpus ends May 2026) and Consensus has not indexed it.
Second-priority abstract retry behind Werderman et al.}
{10.1016/j.atmosenv.2026.122324}

\paperentry{Kang et al.\ 2026}
{Characteristics and dry deposition fluxes of particulate parent, alkyl, and nitro-polycyclic aromatic hydrocarbons in the Bohai Sea Strait, China}
{Atmos Environ}
{\textbf{Metadata only.} Crossref sweep, no abstract deposited. Nitro-PAHs are among the
more directly mutagenic particle-bound species and are under-measured relative to the
parent compounds, and a marine-strait dry-deposition budget is a receptor-side
measurement rather than an emission inventory --- but no flux, no concentration and no
sampling period are extractable from the deposited record. Flagged for retry, low
priority relative to the wind-correction paper.}
{10.1016/j.atmosenv.2026.122322}

\band{Mechanistic toxicology\hfill 2 records}{Violet}

\paperentry{Li \& Li 2026}
{Research advances in the pulmonary carcinogenic effects of atmospheric fine particulate matter mediated by GPER}
{Zhongguo Fei Ai Za Zhi}
{Narrative review positioning the G protein-coupled estrogen receptor as the molecular
hub linking PM$_{2.5}$ to lung carcinogenesis. The argument: environmental estrogen-like
organics and metal ions carried on PM$_{2.5}$ activate GPER, which is overexpressed in
lung tumour tissue, and the activated receptor drives \textbf{MAPK/ERK and PI3K/AKT}
signalling into inflammation, oxidative stress and ferroptosis, with downstream
proliferation, EMT and tumour-microenvironment remodelling. GPER inhibition is proposed
as a target. \textbf{The weaknesses are structural to the format}: no search protocol,
no risk-of-bias assessment, no quantitative synthesis, and a receptor-centric frame that
cannot by construction weigh GPER against the AhR and oxidative-stress pathways it does
not review. Read it as a hypothesis map, not as evidence of pathway dominance. The one
genuinely actionable implication for exposure science is that it assigns the causal role
to specific \emph{components} --- estrogenic organics and transition metals --- rather
than to mass, which is another argument for composition-resolved rather than
gravimetric-equivalent monitoring.}
{10.3779/j.issn.1009-3419.2026.101.13}

\paperentry{Qin et al.\ 2026}
{AlCl$_3$ aerosol exposure is associated with pulmonary lactate accumulation, neutrophil dysfunction, and fibrosis-like remodeling}
{Biofactors}
{Repeated AlCl$_3$ aerosol exposure in rats, with and without the MCT1 inhibitor AZD3965,
plus bone-marrow-derived neutrophils treated with sodium L-lactate, plus a supportive
human cohort of aluminium-exposure-related ILD, non-aluminium ILD and healthy controls
characterised by quantitative CT, spirometry and blood transcriptomics. Exposure produced
pulmonary aluminium deposition, \textbf{lactate accumulation}, oxidative stress, DNA
damage response signalling and collagen-rich fibrosis-like remodelling; in vitro,
high lactate impaired neutrophil chemotaxis, oxidative burst, mitochondrial membrane
potential, calcium signalling and NET responses. AZD3965 partially reversed both.
Al-ILD patients showed greater CT parenchymal abnormality and lower FVC and DLCO than
non-Al-ILD, with innate-immune, cytokine, ECM and NET pathways enriched in blood.
\textbf{No exposure concentration, particle size distribution or duration appears in the
deposited abstract}, so the animal arm cannot be placed on any dose axis, and the human
arm is cross-sectional with no exposure measurement at all --- it supports the mechanism
narratively and cannot test it. Relevant here as occupational-aerosol toxicology and as
a reminder that metal-bearing particulate has organ-level effects that mass concentration
does not index.}
{10.1002/biof.70146}

\band{Respiratory \& allergic\hfill 2 records}{Teal}

\paperentry{Zhongguo Fei Ai Za Zhi 2026}
{China Wumeng Mountain lung cancer: evidence chain, clinical and molecular characteristics of regional lung cancer related to coal pollution in the Yunnan--Guizhou Plateau}
{Zhongguo Fei Ai Za Zhi}
{Narrative synthesis assembling an
exposure\,$\rightarrow$\,component\,$\rightarrow$\,biological-effect\,$\rightarrow$\,%
clinical-molecular\,$\rightarrow$\,prevention chain for household smoky-coal lung cancer
in Xuanwei/Fuyuan, searching PubMed, Web of Science, Embase, CNKI and Wanfang to May 2026.
Carcinogenic components named: \textbf{PAHs and methylated PAHs, PM$_{2.5}$/black carbon,
inorganic minerals and metals, and carbonyls}, as a mixture. The clinical phenotype is
distinctive --- predominantly \textbf{never-smoking women}, high adenocarcinoma fraction,
young onset, early metastasis, multifocal disease and family aggregation. So is the
molecular one: total EGFR mutation is high but classical 19del and L858R are
\emph{under}-represented, while \textbf{G719X, S768I, KRAS G12C and ALK/ROS1 fusions are
elevated} relative to non-Wumeng lung cancer; ARHGEF5, MYO18B and MUC16 are proposed
germline susceptibility genes interacting with POT1/TERT/TERF2. Sputum microbiome shifts
are reported alongside. \textbf{It is a narrative review by an unnamed author group with
no formal inclusion criteria or bias assessment}, drawing on the same group's prior work,
so the evidence chain is an argument rather than a synthesis. Its value is as the clearest
statement in this corpus of a testable source-composition-to-mutation-spectrum hypothesis.}
{10.3779/j.issn.1009-3419.2026.102.16}

\paperentry{Aulia Anwar et al.\ 2026}
{Smoking, environmental exposures, and chronic lung disease as key risk factors for lung cancer in North Sumatra}
{Exp Oncol}
{Hospital-based case-control across three referral hospitals in North Sumatra,
August--November 2024: \textbf{240 histologically confirmed lung cancers and 240 controls}
without malignancy or chronic respiratory disease, matched on sex and residential area
with age adjusted in the model. Independent predictors: \textbf{unhealthy indoor
environment (poor ventilation, household pollutants) OR 2.53 (95\,\% CI 1.45--4.40)} and
\textbf{occupational exposure OR 1.91 (1.13--3.25)}; tuberculosis and COPD were more
prevalent among cases. \textbf{Three weaknesses threaten the inference and should be
stated first.} The exposure is a self-reported composite category with no measured
particulate concentration, so it is subject to recall bias in a design where cases know
their diagnosis. Controls were selected to be free of chronic respiratory disease, which
is a collider on the pathway from household smoke to both COPD and lung cancer and biases
the household-environment OR upward. And smoking differed between groups at $p = 0.023$
yet does not appear among the adjusted independent predictors --- the deposited abstract
does not explain that discrepancy. Treat 2.53 as an upper bound on a poorly specified
exposure, not as a household-PM effect estimate.}
{10.15407/exp-oncology.2026.02.152}

\vspace{4mm}
% ==================================================================== TABLE
\section*{\sffamily\bfseries\color{Deep}Corpus register}
\vspace{-1.5mm}{\color{Amber}\rule{\linewidth}{1.1pt}}\vspace{2.5mm}

{\fontsize{6.7}{8.4}\selectfont
\setlength{\tabcolsep}{2.4pt}
\begin{longtable}{@{}>{\RaggedRight}p{25mm}>{\RaggedRight}p{41mm}>{\RaggedRight}p{27mm}%
>{\RaggedRight}p{23mm}>{\RaggedRight}p{18mm}>{\RaggedRight}p{15mm}@{}}
\toprule
\textsf{\textbf{First author}} & \textsf{\textbf{Focus}} & \textsf{\textbf{Journal}} &
\textsf{\textbf{Design}} & \textsf{\textbf{Metric}} & \textsf{\textbf{Region}}\\
\midrule
\endfirsthead
\toprule
\textsf{\textbf{First author}} & \textsf{\textbf{Focus}} & \textsf{\textbf{Journal}} &
\textsf{\textbf{Design}} & \textsf{\textbf{Metric}} & \textsf{\textbf{Region}}\\
\midrule
\endhead
\multicolumn{6}{@{}l@{}}{\cellcolor{Mist}\textsf{\textbf{\textcolor{Deep}{Neuro / mental health}}}}\\[0.6mm]
\href{https://doi.org/10.1001/jamanetworkopen.2026.25015}{Dales} & Wildfire-sourced PM$_{2.5}$ and ED visits for migraine and other primary headache & \textit{JAMA Netw Open} & Case-crossover & PM$_{2.5}$ (wildfire) & Canada\\
\href{https://doi.org/10.1093/ageing/afag215}{Glover} & Umbrella meta-review of urban environmental exposures and cognitive health & \textit{Age Ageing} & Umbrella review / meta-analysis & PM$_{2.5}$ & Global\\
\href{https://doi.org/10.1016/j.nsa.2026.107019}{Goossens} & In utero air pollution exposure and pre-/postnatal brain development & \textit{Neurosci Appl} & Narrative review & PM$_{2.5}$ / PM$_{10}$ / PM$_{\mathrm{coarse}}$ & Global\\
\href{https://doi.org/10.1016/j.isci.2026.116617}{Ma} & PM$_{2.5}$-induced neurotoxicity and depression: exposure risks and mechanisms & \textit{iScience} & Narrative review & PM$_{2.5}$ + components & Global\\
\href{https://doi.org/10.3760/cma.j.cn112150-20250804-00752}{Pan} & Maternal PM exposure and offspring neurodevelopment: effects and mechanisms & \textit{Zhonghua Yu Fang Yi Xue Za Zhi} & Narrative review & PM (bulk) & China\\
\addlinespace[0.9mm]
\multicolumn{6}{@{}l@{}}{\cellcolor{Mist}\textsf{\textbf{\textcolor{Teal}{Cardiovascular \& metabolic}}}}\\[0.6mm]
\href{https://doi.org/10.3760/cma.j.cn112338-20251120-00843}{Huang} & PM$_{2.5}$ and dyslipidaemia in rural Chinese children; dietary effect modification & \textit{Zhonghua Liu Xing Bing Xue Za Zhi} & Cross-sectional (multistage) & PM$_{2.5}$ & China\\
\href{https://doi.org/10.32394/pe/218966}{Kazmierska} & Air pollution and cardiovascular disease: review of current data & \textit{Przegl Epidemiol} & Narrative review & PM$_{2.5}$ / PM$_{10}$ & Global\\
\href{https://doi.org/10.1007/s00125-026-06814-2}{Xi} & Air pollution, built environment and temperature vs. prediabetes phenotypes & \textit{Diabetologia} & Prospective cohort & PM$_{2.5}$ & Germany\\
\href{https://doi.org/10.1161/JAHA.125.046917}{Yang} & Short-term PM$_{2.5}$ and hospitalisation/ED use in atrial fibrillation & \textit{J Am Heart Assoc} & Case-crossover & PM$_{2.5}$ (wildfire-influenced) & USA\\
\addlinespace[0.9mm]
\multicolumn{6}{@{}l@{}}{\cellcolor{Mist}\textsf{\textbf{\textcolor{Sky}{Reproductive \& developmental}}}}\\[0.6mm]
\href{https://doi.org/10.1016/j.jes.2025.10.043}{Chen} & Prenatal air pollution, fetal growth trajectory phenotypes and child neurodevelopment & \textit{J Environ Sci (China)} & Prospective cohort & PM$_{2.5}$ / PM$_{10}$ & China\\
\href{https://doi.org/10.1016/j.ecoenv.2026.120566}{Huang} & Joint prenatal bisphenol + air-pollutant mixture exposure and preschool IQ & \textit{Ecotoxicol Environ Saf} & Prospective cohort (mixtures) & PM$_{2.5}$ / PM$_{10}$ & China\\
\href{https://doi.org/10.1007/s40572-026-00552-8}{Peterson} & Ambient air pollution, PAHs and pubertal development: critical review & \textit{Curr Environ Health Rep} & Critical review & PM$_{2.5}$ components & Global\\
\href{https://doi.org/10.1016/j.envres.2026.125313}{Xie} & Long-term PM$_{2.5}$ and timing of pubertal development in boys and girls & \textit{Environ Res} & Prospective cohort & PM$_{2.5}$ & China\\
\addlinespace[0.9mm]
\multicolumn{6}{@{}l@{}}{\cellcolor{Mist}\textsf{\textbf{\textcolor{Sage}{Respiratory \& allergic}}}}\\[0.6mm]
\href{https://doi.org/10.1007/s00484-026-03289-z}{Bayar Muluk} & Climate change and otolaryngology: airway disease impacts and emission reduction & \textit{Int J Biometeorol} & Narrative review & PM$_{2.5}$ & Global\\
\href{https://doi.org/10.4103/mgr.MEDGASRES-D-25-00164}{Liao} & Short-term air pollution, respiratory symptoms and lung function in asthma & \textit{Med Gas Res} & Panel study & PM$_{2.5}$ / PM$_{10}$ & China\\
\href{https://doi.org/10.3390/epidemiologia7040097}{Sanchez de Prada} & Ambient PM and asthma-coded admissions in children 0-3 y, Spain & \textit{Epidemiologia} & Ecological panel & PM$_{2.5}$ / PM$_{10}$ & Spain\\
\href{https://doi.org/10.1016/j.jhazmat.2026.143076}{Yang} & Allergenic vs. toxic risk of airborne seed hairs compared with ambient PM$_{2.5}$ & \textit{J Hazard Mater} & Field measurement & PM$_{2.5}$ + bioaerosol & China\\
\addlinespace[0.9mm]
\multicolumn{6}{@{}l@{}}{\cellcolor{Mist}\textsf{\textbf{\textcolor{Violet}{Mechanistic toxicology}}}}\\[0.6mm]
\href{https://doi.org/10.1016/j.jes.2025.12.074}{Ding} & PM$_{2.5}$ promotes vascular calcification via miR-27b-5p/HES1 signalling & \textit{J Environ Sci (China)} & In vitro + organoid + murine & PM$_{2.5}$ & China\\
\href{https://doi.org/10.1007/s00360-026-01698-5}{Maraschi} & Osmoregulatory disruption and metal bioaccumulation from settleable atmospheric PM & \textit{J Comp Physiol B} & Controlled exposure (ecotox) & Settleable APM & Brazil\\
\href{https://doi.org/10.1016/j.jes.2025.11.018}{Wu} & Macrophage NLRP3 inflammasome in real-ambient PM-induced hepatic glucose dysregulation & \textit{J Environ Sci (China)} & Real-ambient murine + co-culture & Real-ambient PM & China\\
\addlinespace[0.9mm]
\multicolumn{6}{@{}l@{}}{\cellcolor{Mist}\textsf{\textbf{\textcolor{Clay}{Occupational \& indoor}}}}\\[0.6mm]
\href{https://doi.org/10.3389/fmedt.2026.1796828}{Chen} & Nested suction instrument to cut intraoperative surgical-smoke PM$_{2.5}$ & \textit{Front Med Technol} & Bench / device experiment & PM$_{2.5}$ (surgical smoke) & China\\
\href{https://doi.org/10.1093/annweh/wxag055}{Cole} & Determinants of respirable crystalline silica exposure in Australian tunnelling & \textit{Ann Work Expo Health} & Observational exposure study & Respirable dust / RCS & Australia\\
\href{https://doi.org/10.1016/j.jhazmat.2026.143081}{Ding} & Stage-resolved hair Zn-isotope dynamics as a firefighting combustion-exposure tracer & \textit{J Hazard Mater} & Repeated-measures biomarker & Combustion-derived PM & China\\
\href{https://doi.org/10.3390/jox16040119}{Ramsden} & Xenobiotic hazards in aircraft cabin air & \textit{J Xenobiot} & Review + risk economics & Ultrafine particles & Global\\
\addlinespace[0.9mm]
\multicolumn{6}{@{}l@{}}{\cellcolor{Mist}\textsf{\textbf{\textcolor{Amber}{Sensing, forecasting \& instrumentation}}}}\\[0.6mm]
\href{https://doi.org/10.1016/j.jenvman.2026.130515}{Shi} & Adaptive spatiotemporal graph transformer for multi-site multi-step PM$_{2.5}$ forecasting & \textit{J Environ Manage} & Deep-learning model & PM$_{2.5}$ & China / India\\
\href{https://doi.org/10.3390/bios16070380}{Tsuchiya} & Seasonal and indoor-outdoor characteristics of fluorescent aerosol particles in offices & \textit{Biosensors} & Real-time sensor campaign & Bioaerosol / coarse & Japan\\
\addlinespace[0.9mm]
\multicolumn{6}{@{}l@{}}{\cellcolor{Mist}\textsf{\textbf{\textcolor{Amber}{Exposure assessment \& modelling}}}}\\[0.6mm]
\href{https://doi.org/10.1029/2025GH001621}{Solorzano-Araque} & Multiscale exposure modelling: LOTOS-EUROS + mobility surveys + particle toxicity & \textit{GeoHealth} & CTM + mobility integration & PM$_{2.5}$ / PM$_{10}$ & Colombia\\
\addlinespace[0.9mm]
\multicolumn{6}{@{}l@{}}{\cellcolor{Mist}\textsf{\textbf{\textcolor{Coral}{Burden, policy \& mitigation}}}}\\[0.6mm]
\href{https://doi.org/10.1002/hsr2.72886}{Karamzad} & Ischaemic heart disease burden in MENA 1990-2021; PM attributable fraction 28.2\% & \textit{Health Sci Rep} & GBD secondary analysis & Ambient PM & MENA\\
\href{https://doi.org/10.1371/journal.pone.0347989}{Rees} & HEPA filtration to reduce respiratory infection in care homes & \textit{PLoS One} & RCT process evaluation & Filtered PM & UK\\
\href{https://doi.org/10.1016/j.jes.2026.01.009}{Shan} & Policy-driven emission reductions and post-pandemic rebound, high-resolution inventories & \textit{J Environ Sci (China)} & Emission inventory & PM emissions & China\\
\addlinespace[0.9mm]
\multicolumn{6}{@{}l@{}}{\cellcolor{Mist}\textsf{\textbf{\textcolor{Slate}{Other clinical endpoints}}}}\\[0.6mm]
\href{https://doi.org/10.3390/dj14070440}{Lupita} & PM$_{2.5}$, preventive-care capacity and caries/edentulism burden across the EU & \textit{Dent J} & Ecological panel & PM$_{2.5}$ & EU-27\\
\href{https://doi.org/10.1016/j.jdin.2026.06.008}{Park} & Long-term PM exposure and risk of androgenetic alopecia & \textit{JAAD Int} & Nationwide cohort & PM$_{2.5}$ / PM$_{10}$ & South Korea\\
\href{https://doi.org/10.3171/2026.3.PEDS2615}{Sun} & Environmental exposure and area deprivation vs. paediatric hydrocephalus outcomes & \textit{J Neurosurg Pediatr} & Registry cohort + geospatial & PM$_{2.5}$ (satellite-derived) & USA\\
\addlinespace[0.9mm]
\bottomrule
\end{longtable}}


\vspace{2mm}
\begin{keybox}
{\sffamily\bfseries\footnotesize\color{Deep}Provenance and method}\\[1.2mm]
{\fontsize{7.2}{8.8}\selectfont
Entry window \textbf{2026-08-24 $\rightarrow$ 2026-08-24}, contiguous with the 23 August
issue. \textbf{This issue was built late}, on 25 August: the 24 August harvest ran and its
cache was warmed, but the issue was deliberately deferred to the following scheduled run.
The window is unchanged and no gap was left. \textbf{Sources.} PubMed E-utilities on
\texttt{[EDAT]} returned 3 records; the connector re-query on the health axis returned
\textbf{4} and on the instrumentation axis \textbf{1}, five distinct records in total,
of which four shipped. Europe PMC on \texttt{CREATION\_DATE} returned \textbf{9 records,
all archival PMC back-scans} with no title, DOI or abstract --- the same structural
artefact as the 13 seen on 23 August --- and none shipped. The Crossref by-ISSN sweep
returned 9 records and shipped 5, including the \textbf{first four \textit{Atmospheric
Chemistry and Physics} deposits the pipeline has ever caught}. \textbf{OpenAlex returned
an empty result for the 22nd consecutive run} and remains a broken leg. arXiv returned
nothing new under the recency screen. Consensus was swept on low-cost sensor calibration
and drift and on satellite--surface PM$_{2.5}$ fusion and returned 6 on-topic records,
\textbf{all 6 already in \texttt{seen.json}}. Scholar Gateway was queried for the
abstract-free WRF--CMAQ record and cannot reach it; its corpus ends May 2026.
\textbf{Screening.} 21 records were screened out and written to
\texttt{state/rejected.jsonl} with an individual reason each; the 9 archival back-scans
carry a shared structural reason and the 6 Consensus hits carry the issue date on which
each was originally shipped. One borderline exclusion is logged as borderline: a
physics-informed indoor CO$_2$ prediction method whose architecture transfers directly to
PM but which reports no particulate contrast. \textbf{Labels.} Two records carry no
abstract and are marked \textit{Metadata only}; their design, sample size, endpoint and
study location are recorded as not extracted rather than inferred. \textbf{The forest plot
carries two estimates from a single model} and both exposures are self-reported
categories, stated on the axis; no pooled estimate is drawn. All DOIs were checked by
\texttt{check\_dois.py} before publication. Content derived from publisher metadata;
abstracts remain the copyright of their respective publishers.\par}
\end{keybox}

\end{document}
